% =====================================================================
%  BAB 7 — KESIMPULAN DAN SARAN
% =====================================================================

\chapter{Kesimpulan dan Saran}

\section{Kesimpulan}
\label{sec:7.1}

Penelitian ini telah berhasil [menjawab/merancang/membangun/membuktikan] [tujuan utama
penelitian]. Berdasarkan hasil pengujian dan analisis yang telah dilakukan, dapat
ditarik kesimpulan sebagai berikut.

\begin{enumerate}
  \item [Kesimpulan 1 — jawaban langsung atas Rumusan Masalah 1. Sertakan angka/fakta
        konkret dari hasil pengujian, bukan pernyataan umum.]

  \item [Kesimpulan 2 — jawaban atas Rumusan Masalah 2. Spesifik dan terukur.]

  \item [Kesimpulan 3 — jawaban atas Rumusan Masalah 3. Bandingkan dengan baseline atau
        penelitian terdahulu jika relevan.]
\end{enumerate}

[Paragraf penutup yang merangkum kontribusi utama penelitian dan signifikansinya terhadap
bidang keilmuan yang bersangkutan.]

\section{Saran}
\label{sec:7.2}

Berdasarkan temuan dan keterbatasan yang diidentifikasi, beberapa saran untuk penelitian
dan pengembangan lebih lanjut adalah sebagai berikut.

\begin{enumerate}
  \item \textbf{[Saran 1: Perbaikan teknis yang spesifik.]}
        [Jelaskan dengan singkat mengapa ini penting dan bagaimana bisa dilakukan.]

  \item \textbf{[Saran 2: Perluasan cakupan eksperimen.]}
        [Misalnya: menguji pada dataset yang lebih beragam, lingkungan produksi nyata,
        atau skenario serangan yang lebih kompleks.]

  \item \textbf{[Saran 3: Arah penelitian lanjutan.]}
        [Pertanyaan terbuka yang muncul dari penelitian ini yang layak dieksplorasi
        lebih jauh.]

  \item \textbf{[Saran 4: Aspek yang belum dicakup dalam penelitian ini.]}
        [Misalnya: integrasi dengan sistem yang lebih besar, evaluasi biaya operasional,
        atau aspek privasi dan etika yang relevan.]
\end{enumerate}
